%! Author = tstreule

\section{Deep Gradients}\label{sec:deep-gradients}


% ===
\subsection{Short connectivity}

\emph{Residual connections:}\enskip
shortcut layers and \textbf{add} back in,\enskip $g(\symbf x) = f(\symbf{x; \theta}) \green + \symbf x$, or if dimensions don't match: $f(\symbf{x; \theta}) \green + \symbf{Wx}$

\emph{Skip connections:}\enskip
shortcut layers and \textbf{concatenate} back in.
Connect layer output to all downstream layers.


% ===
\subsection{Normalization}

We observe \textbf{vanishing gradients}, i.e. poor backpropagation of errors, when incoming data to activation function is not normalized.


\subsubsection{Batch normalization \enskip\mdseries $\implies$ across many data points}

\begin{itemize}
    \item Learning rate $\eta$ is otherwise very sensitive
    \item Key idea: normalize the layer activation (and backpropagate through normalization!)
\end{itemize}
($+$) very effective, broadly used in vision\\
($-$) dependence on batch size, not suitable for all architectures


\subsubsection{Layer normalization \enskip\mdseries $\implies$ across units in same layer}

\begin{itemize}
    \item In scenarios like recurrent NNs, we have batch size = 1; batch normalization is not possible.
    \item Use population average of activations in layer as reference
\end{itemize}
($+$) very effective, used in NLP (transformers, RNNs)\\
($-$) no consistent benefits
